% ============================================================================
%  Załączniki
% ============================================================================
\chapter*{Załączniki}
\addcontentsline{toc}{chapter}{Załączniki}

\section*{Załącznik A. Kod źródłowy i materiały dodatkowe}
\addcontentsline{toc}{section}{Załącznik A. Kod źródłowy i materiały dodatkowe}

Pełny kod źródłowy gry wraz z~projektem Unity (sceny, prefabrykaty, materiały, pliki
konfiguracyjne) znajduje się w~dołączonym do pracy katalogu z~kodem źródłowym oraz
w~repozytorium \url{https://github.com/neavailable/Praca-Nazarii-Yasinskyi} (kod gry:
\kod{Assets/Code}, skrypty edytora: \kod{Assets/Editor}). Katalog \kod{praca\_dyplomowa}
zawiera źródła pracy w~formacie \LaTeX, zrzuty ekranu oraz skrypt
\sciezka{narzedzia/ThesisCapture.cs} realizujący scenariusz automatyczny opisany
w~rozdziale~\ref{rozdz:testy}, przeznaczony do uruchamiania na kopii projektu.

\section*{Załącznik B. Parametry rozgrywki zapisane w plikach konfiguracyjnych}
\addcontentsline{toc}{section}{Załącznik B. Parametry rozgrywki zapisane w plikach konfiguracyjnych}

Tabela zawiera wszystkie wartości zapisane w~plikach JSON w~katalogu
\sciezka{Assets/Resources/Configurations}. Nagłówki grup podają nazwę pliku (bez przyrostka
\kod{Configuration.json}) i~sekcję; każdej z~25 sekcji odpowiada jedna klasa konfiguracji.

\newcommand{\grupa}[2]{\multicolumn{2}{@{}l}{\rule{0pt}{9pt}\textit{#1 -- #2}}\\}
\begingroup
\footnotesize
\setstretch{1.0}
\renewcommand{\arraystretch}{0.98}
\noindent
\begin{minipage}[t]{0.49\textwidth}
\begin{tabular}[t]{@{}p{5.7cm}r@{}}
\toprule
\textbf{Parametr} & \textbf{Wartość}\\
\midrule
\grupa{Ball}{Movement}
\hspace{0.3cm}\texttt{MaxPushingPower}, \texttt{PushPowerDecrease} & 25,0; 20,0\\
\grupa{Ball}{BounceAnimation}
\hspace{0.3cm}\texttt{AnimationDuration} & 0,2\\
\hspace{0.3cm}\texttt{SqueezingExtendingSpeedDivider} & 1,8\\
\hspace{0.3cm}\texttt{ColorChangeSpeed} & 6,0\\
\multicolumn{2}{@{}l}{\hspace{0.3cm}\texttt{TargetColorR/G/B/A}: 218/67/67/255}\\
\grupa{Ball Menu}{Movement}
\hspace{0.3cm}\texttt{StartPushingPower} & 7,0\\
\grupa{Boost Pad}{Boost Pad}
\hspace{0.3cm}\texttt{BoostPower}, \texttt{MaxSpeed} & 20,0; 35,0\\
\grupa{Cactus}{Movement / Damage}
\hspace{0.3cm}\texttt{Speed}, \texttt{Damage} & 3; 29\\
\grupa{Coin}{Value}
\hspace{0.3cm}\texttt{CoinValue} & 20\\
\grupa{Coin}{Movement / Rotation}
\hspace{0.3cm}\texttt{YSpeed}, \texttt{YBorder} & 0,15; 0,2\\
\hspace{0.3cm}\texttt{YRotation} & 2,0\\
\grupa{Coin}{Destroyer}
\hspace{0.3cm}\texttt{ShrinkingSpeed}, \texttt{MinScale} & 2,0; 0,05\\
\grupa{Dust Effect}{Position}
\hspace{0.3cm}\texttt{YOffset} & 0,2\\
\grupa{Enemy}{Bounce}
\hspace{0.3cm}\texttt{LostPower} & 0,2\\
\bottomrule
\end{tabular}
\end{minipage}\hfill
\begin{minipage}[t]{0.49\textwidth}
\begin{tabular}[t]{@{}p{5.6cm}r@{}}
\toprule
\textbf{Parametr} & \textbf{Wartość}\\
\midrule
\grupa{Ground}{Friction}
\hspace{0.3cm}\texttt{FrictionPower} & 4,0\\
\grupa{Heal}{Heal}
\hspace{0.3cm}\texttt{HealAmount} & 30\\
\grupa{Loading}{LoadingTime}
\hspace{0.3cm}\texttt{MinLoadingTime}, \texttt{MaxLoadingTime} & 1; 4\\
\grupa{Minimap}{Minimap}
\hspace{0.3cm}\texttt{OrthographicSize}, \texttt{Height} & 14,0; 40,0\\
\hspace{0.3cm}\texttt{TextureSize} & 256\\
\grupa{Mushroom}{Rotation / Damage}
\hspace{0.3cm}\texttt{MinAngle}, \texttt{MaxAngle} & 0, 360\\
\hspace{0.3cm}\texttt{Damage} & 17\\
\grupa{Obstacles}{Obstacles}
\hspace{0.3cm}\texttt{BumperPower} & 12,0\\
\hspace{0.3cm}\texttt{BumperMaxSpeed} & 30,0\\
\hspace{0.3cm}\texttt{KnockPower} & 10,0\\
\hspace{0.3cm}\texttt{SpinnerSpeed} & 90,0\\
\hspace{0.3cm}\texttt{MovingWallSpeed} & 2,5\\
\hspace{0.3cm}\texttt{MovingWallDistance} & 2,5\\
\grupa{Push}{Push}
\hspace{0.3cm}\texttt{PushPower} & 8,0\\
\grupa{UI}{Coin}
\hspace{0.3cm}\texttt{ChangeSpeed}, \texttt{ScaleFontSpeed} & 0,05; 100\\
\hspace{0.3cm}\texttt{MaxFontSize} & 33\\
\grupa{UI}{Health}
\hspace{0.3cm}\texttt{Delay}, \texttt{DefaultHealth} & 0,05; 100\\
\hspace{0.3cm}\texttt{MaxFontSize} & 40\\
\grupa{Wall / Wall Menu}{Bounce}
\hspace{0.3cm}\texttt{LostPower} & 0,7 / 1,0\\
\grupa{Wind}{Wind}
\hspace{0.3cm}\texttt{MinStrength}, \texttt{MaxStrength} & 0,3; 0,8\\
\bottomrule
\end{tabular}
\end{minipage}
\endgroup

\section*{Załącznik C. Instrukcja uruchomienia i budowania aplikacji}
\addcontentsline{toc}{section}{Załącznik C. Instrukcja uruchomienia i budowania aplikacji}

Wymagany jest edytor Unity 2022.3.60f1, a~do budowania wersji Android także moduły
Android Build Support, Android SDK \& NDK Tools oraz OpenJDK, instalowane w~programie
Unity Hub. Pierwsze otwarcie projektu i~pierwsze budowanie wymagają połączenia
z~internetem (pakiet UniTask oraz zależności Androida wtyczek Firebase i~Appodeal).
\begin{enumerate}[leftmargin=1.9cm, itemsep=0pt]
  \item Otworzyć katalog projektu w~programie Unity Hub i~poczekać na import zasobów.
  \item Aby zagrać w~edytorze, otworzyć scenę \sciezka{Assets/Scenes/Bootstrap.unity}
        i~nacisnąć przycisk odtwarzania; strzał wykonuje się, przeciągając myszą od kuli.
  \item Aby zbudować plik APK, wybrać polecenie \emph{Build $\rightarrow$ Android APK}
        (plik APK trafia do katalogu \sciezka{Builds/Android}) lub uruchomić skrypt
        \sciezka{./Build-Android.ps1}, który buduje grę w~osobnej kopii projektu.
  \item Parametry rozgrywki zmienia się w~plikach JSON w~katalogu
        \sciezka{Assets/Resources/Configurations}; są one sprawdzane przed każdym budowaniem.
\end{enumerate}
